\textbf{Phần 1. Trắc nghiệm lựa chọn}
\setcounterpageref{ex}{0}

\begin{ex} %[2D2V08]
Kết quả của phép tính $9 + 7$ là
\choice
{$15$}
{\True $16$}
{$17$}
{$18$}
\loigiai{
$9 + 7 = 16$.
\textbf{Đáp án của câu hỏi là B}
}
\end{ex}

\begin{ex} %[2D2V08]
Trong vườn có $8$ cây cam và số cây bưởi bằng số cây cam cộng thêm $5$ cây. Hỏi trong vườn có tất cả bao nhiêu cây cam và bưởi?
\choice
{$13$ cây}
{$18$ cây}
{$19$ cây}
{\True $21$ cây}
\loigiai{
Số cây bưởi là: $8 + 5 = 13$ (cây).\nTổng số cây cam và bưởi là: $8 + 13 = 21$ (cây).
\textbf{Đáp án của câu hỏi là D}
}
\end{ex}

\begin{ex} %[2D2V08]
Tính nhanh giá trị: $8 + 5 + 2 + 5$.
\choice
{\True $20$}
{$18$}
{$19$}
{$21$}
\loigiai{
$(8 + 2) + (5 + 5) = 10 + 10 = 20$.
\textbf{Đáp án của câu hỏi là A}
}
\end{ex}

\begin{ex} %[2D2V08]
Tìm một số biết rằng lấy số đó trừ đi $9$ thì được số chẵn nhỏ nhất có hai chữ số.
\choice
{$18$}
{$20$}
{\True $19$}
{$17$}
\loigiai{
Số chẵn nhỏ nhất có hai chữ số là 10. Số đó là: $10 + 9 = 19$.
\textbf{Đáp án của câu hỏi là C}
}
\end{ex}

\begin{ex} %[2D2V08]
Hộp thứ nhất có $7$ chiếc bút chì, hộp thứ hai có $8$ chiếc bút chì. Bạn Minh lấy ra $4$ chiếc bút chì từ hộp thứ nhất. Hỏi cả hai hộp còn lại bao nhiêu chiếc bút chì?
\choice
{$12$ chiếc}
{\True $11$ chiếc}
{$10$ chiếc}
{$13$ chiếc}
\loigiai{
Tổng số bút lúc đầu: $7 + 8 = 15$ (chiếc).\nSau khi lấy ra còn: $15 - 4 = 11$ (chiếc).
\textbf{Đáp án của câu hỏi là B}
}
\end{ex}

\begin{ex} %[2D2V08]
Điền dấu thích hợp: $6 + 8 \dots 7 + 7$.
\choice
{$>$}
{$<$}
{Không so sánh được}
{\True $=$}
\loigiai{
$6 + 8 = 14$; $7 + 7 = 14$. Hai vế bằng nhau nên điền dấu $=$.
\textbf{Đáp án của câu hỏi là D}
}
\end{ex}

\begin{ex} %[2D2V08]
Một sợi dây dài $18\text{ cm}$. Người ta cắt thành hai đoạn, biết đoạn thứ nhất dài $9\text{ cm}$. Đoạn thứ hai dài bao nhiêu?
\choice
{\True $9\text{ cm}$}
{$8\text{ cm}$}
{$10\text{ cm}$}
{$7\text{ cm}$}
\loigiai{
Đoạn thứ hai dài: $18 - 9 = 9\text{ (cm)}$.
\textbf{Đáp án của câu hỏi là A}
}
\end{ex}

\begin{ex} %[2D2V08]
Tìm số tự nhiên $y$ biết: $y + 7 = 16$.
\choice
{$8$}
{$10$}
{\True $9$}
{$7$}
\loigiai{
$y = 16 - 7 = 9$.
\textbf{Đáp án của câu hỏi là C}
}
\end{ex}

\begin{ex} %[2D2V08]
Bố mua cho Mai $8$ quyển vở, mẹ mua cho Mai thêm $6$ quyển vở nữa. Mai đã dùng hết $5$ quyển. Hỏi Mai còn lại bao nhiêu quyển vở chưa dùng?
\choice
{$8$ quyển}
{\True $9$ quyển}
{$10$ quyển}
{$7$ quyển}
\loigiai{
Tổng số vở Mai có: $8 + 6 = 14$ (quyển).\nSố vở còn lại: $14 - 5 = 9$ (quyển).
\textbf{Đáp án của câu hỏi là B}
}
\end{ex}

\begin{ex} %[2D2V08]
Tổng của hai số bằng $15$. Nếu tăng số hạng thứ nhất thêm $2$ đơn vị và tăng số hạng thứ hai thêm $3$ đơn vị thì tổng mới là
\choice
{$18$}
{$19$}
{$21$}
{\True $20$}
\loigiai{
Tổng mới là: $15 + 2 + 3 = 20$.
\textbf{Đáp án của câu hỏi là D}
}
\end{ex}

\begin{ex} %[2D2V08]
Hai số có hiệu bằng $8$, số trừ là $9$. Số bị trừ là
\choice
{\True $17$}
{$16$}
{$18$}
{$15$}
\loigiai{
Số bị trừ = Hiệu + Số trừ = $8 + 9 = 17$.
\textbf{Đáp án của câu hỏi là A}
}
\end{ex}

\begin{ex} %[2D2V08]
Mai có $9$ bông hoa, Cúc có $7$ bông hoa, Lan có số hoa bằng số hoa của Mai và Cúc cộng lại rồi bớt đi $4$ bông. Lan có bao nhiêu bông hoa?
\choice
{$11$ bông hoa}
{$13$ bông hoa}
{\True $12$ bông hoa}
{$10$ bông hoa}
\loigiai{
Số hoa của Lan: $9 + 7 - 4 = 16 - 4 = 12$ (bông hoa).
\textbf{Đáp án của câu hỏi là C}
}
\end{ex}

\textbf{Phần 2. Trắc nghiệm đúng sai}
\setcounterpageref{ex}{0}

\begin{ex} %[2D2V08]
Xét tính đúng sai của các phép tính trong bảng cộng có nhớ:
\choiceTF
{\True Phép cộng $6 + 7$ có kết quả bằng $13$.}
{Tổng của số lẻ lớn nhất có 1 chữ số và số chẵn nhỏ nhất có 1 chữ số khác 0 là 10.}
{\True Biểu thức $8 + 9$ có giá trị lớn hơn biểu thức $7 + 9$.}
{\True Nếu $a + 6 = 15$ thì $a = 9$.}
\loigiai{
\begin{itemize}
    \item Ý a) Đúng: $6 + 7 = 13$.
    \item Ý b) Sai: Số lẻ lớn nhất có 1 chữ số là 9, số chẵn nhỏ nhất khác 0 là 2. Tổng là $9 + 2 = 11$ (không phải 10).
    \item Ý c) Đúng: $8 + 9 = 17 > 7 + 9 = 16$.
    \item Ý d) Đúng: $15 - 6 = 9$. Đúng.
\end{itemize}
\begin{itemize}
    \item a) \textbf{Đúng}.
    \item b) \textbf{Sai}.
    \item c) \textbf{Đúng}.
    \item d) \textbf{Đúng}.
\end{itemize}
}
\end{ex}

\begin{ex} %[2D2V08]
Một đàn gà có 9 con gà trống và 8 con gà mái:
\choiceTF
{\True Tổng số gà trong đàn là 17 con.}
{Số gà mái nhiều hơn số gà trống 1 con.}
{\True Nếu mua thêm 3 con gà mái thì số gà mái sẽ nhiều hơn số gà trống 2 con.}
{\True Tổng số chân của 1 con gà trống và 1 con gà mái là 4 chân.}
\loigiai{
\begin{itemize}
    \item Ý a) Đúng: $9 + 8 = 17$ con.
    \item Ý b) Sai: Gà trống có 9 con, mái có 8 con nên gà trống nhiều hơn gà mái 1 con.
    \item Ý c) Đúng: Gà mái lúc này là $8 + 3 = 11$. $11 - 9 = 2$ con. Đúng.
    \item Ý d) Đúng: Mỗi con gà có 2 chân, $2 + 2 = 4$ chân.
\end{itemize}
\begin{itemize}
    \item a) \textbf{Đúng}.
    \item b) \textbf{Sai}.
    \item c) \textbf{Đúng}.
    \item d) \textbf{Đúng}.
\end{itemize}
}
\end{ex}

\textbf{Phần 3. Trả lời ngắn}
\setcounterpageref{ex}{0}

\begin{ex} %[2D2V08]
Tính giá trị biểu thức: $9 + 8 - 7$.
\shortans{10}
\loigiai{
\textbf{Đáp số:} $10$\\
$9 + 8 = 17$; $17 - 7 = 10$.
}
\end{ex}

\begin{ex} %[2D2V08]
Nam có 8 viên bi, Hoàng có nhiều hơn Nam 5 viên bi. Hỏi cả hai bạn có tất cả bao nhiêu viên bi?
\shortans{21}
\loigiai{
\textbf{Đáp số:} $21$\\
Hoàng có: $8 + 5 = 13$ viên. Cả hai bạn: $8 + 13 = 21$ viên.
}
\end{ex}

\begin{ex} %[2D2V08]
Tìm số tự nhiên a biết: $a + 8 = 17$.
\shortans{9}
\loigiai{
\textbf{Đáp số:} $9$\\
$a = 17 - 8 = 9$.
}
\end{ex}

\begin{ex} %[2D2V08]
Tìm tổng của số lớn nhất có 1 chữ số và số liền sau của nó.
\shortans{19}
\loigiai{
\textbf{Đáp số:} $19$\\
Số lớn nhất có 1 chữ số là 9. Số liền sau là 10. Tổng: $9 + 10 = 19$.
}
\end{ex}

\begin{ex} %[2D2V08]
Lớp 2A có 9 bạn tham gia thi hát, số bạn thi vẽ nhiều hơn thi hát 4 bạn. Hỏi có tất cả bao nhiêu bạn tham gia thi cả hai môn?
\shortans{22}
\loigiai{
\textbf{Đáp số:} $22$\\
Thi vẽ: $9 + 4 = 13$. Tổng số: $9 + 13 = 22$ bạn.
}
\end{ex}

\begin{ex} %[2D2V08]
Có bao nhiêu số có một chữ số lớn hơn 6 mà khi cộng với 9 cho kết quả là số có hai chữ số?
\shortans{3}
\loigiai{
\textbf{Đáp số:} $3$\\
Các số lớn hơn 6 là 7, 8, 9 (có 3 số).
}
\end{ex}
